\documentclass[11pt]{article}
\usepackage[margin=25mm]{geometry}
\usepackage{amsmath,amssymb}
\begin{document}

\title{I01 prevalidation 15 protocol:\\Dimension scan of the paired-transfer selection rule}
\date{2026-10-04}
\maketitle

\section*{Question}
Stage 14 established a four-class support-transfer pattern in the $n=3$ energy
correlated clock circuit. This protocol tests whether the same support classes remain
distinguishable when the number of momentum pairs is changed.

\section*{Frozen exact-Aer conditions}
For $n=2,3,4$, prepare the $P=0$ thermal state with the same cutoff, beta, reference
weight, transverse $R_x(+1)$--$R_x(-1)$ pair, skewed $2n+1$ phase labels, two blocks,
and $\tau=0$. The reference clock uses enough bits to encode base energies through
$n$. At each dimension evaluate four fixed source-return classes: $(0,0)$, $(0,1)$,
$(0,2)$, and $(0,3)$. All science axes are evaluated explicitly.

\section*{Validation}
For $n=3$, the generalized circuit must reproduce the corresponding stage-14 matrix
entries before its $n=2$ and $n=4$ rows are interpreted. Clock acceptance is checked
at every circuit setting and must remain $2^{-c}$ for $c$ clock qubits.

\section*{Scope}
This is a finite-family toy-model robustness check. It does not claim that $n=3$ is
physically preferred, nor that an observed rank constitutes a spontaneous dimension
selection mechanism.

\section*{Reproduction}
\texttt{../src/relational\_atlas\_i01/holdout/qiskit\_rank\_holdout15\_dimension\_transfer.py}
\end{document}
